\section{Benutzertests}

\subsection{Usability-Aufgaben}

Basierend auf der Aufgabenanalyse aus Meilenstein~1
wurden vier Usability-Aufgaben entwickelt, die typische Nutzungsszenarien
der App abdecken.
Die Testpersonen wurden gebeten, laut zu denken (Think-Aloud-Methode).

\subsubsection*{Aufgabe 1 – Kartennavigation und Pin-Inspektion}

\textbf{Szenario:}
\textit{„Sie sind gerade in Wien unterwegs und möchten spontan mehr über
das Kunsthistorische Museum erfahren.
Finden Sie es auf der Karte, tippen Sie darauf und lesen Sie die
angezeigten Informationen."}

\textbf{Ziel:} Testen, ob die farbkodierten Pins auffindbar und antippbar sind
und ob das Detail-Bottom-Sheet intuitiv geöffnet und gelesen werden kann.

\subsubsection*{Aufgabe 2 – Datensatz suchen und herunterladen}

\textbf{Szenario:}
\textit{„Sie planen eine Radtour durch Wien und möchten offline auf
Radwegdaten zugreifen können.
Finden Sie im Downloads-Bereich den Datensatz \emph{Radwege Wien}
und laden Sie ihn herunter."}

\textbf{Ziel:} Testen, ob die Suchfunktion im Downloads-Tab und
der Download-Flow klar und verständlich sind.

\subsubsection*{Aufgabe 3 – Suche auf der Karte}

\textbf{Szenario:}
\textit{„Sie möchten schnell den Standort von \emph{Schloss Schönbrunn}
auf der Karte sehen.
Verwenden Sie die Suchfunktion der App, um Schönbrunn zu finden
und auf der Karte anzuzeigen."}

\textbf{Ziel:} Testen, ob das Such-Bottom-Sheet auf der Kartenansicht
auffindbar ist und ob die Kamera korrekt auf das Suchergebnis zoomt.

\subsubsection*{Aufgabe 4 – Einstellungen anpassen und Datensatz teilen}

\textbf{Szenario:}
\textit{„Aktivieren Sie den Dunkelmodus in den Einstellungen.
Teilen Sie anschließend den Datensatz \emph{Wiener Baumkataster}
mit jemandem."}

\textbf{Ziel:} Testen, ob die Einstellungen ohne Erklärung gefunden werden
und ob die Swipe-to-reveal-Geste für die Share-Funktion entdeckt wird.

\subsection{Fragenkatalog}

\subsubsection*{Fragen nach jeder Aufgabe (Post-Task)}

\begin{enumerate}
  \item Wie schwierig war diese Aufgabe für Sie?\\
        \textit{(Skala: 1 = sehr einfach \ldots\ 7 = sehr schwierig)}
  \item Gab es einen Moment, in dem Sie nicht wussten, wie Sie
        weiter vorgehen sollen? Falls ja: Wo genau?
  \item Was haben Sie erwartet zu passieren, als Sie auf
        [jeweilige Schaltfläche / Bereich] getippt haben?
  \item Was würden Sie an dieser Funktion ändern?
\end{enumerate}

\subsubsection*{System Usability Scale (SUS) – nach allen Aufgaben}

Die folgenden zehn Aussagen werden auf einer Skala von 1
(\textit{stimme überhaupt nicht zu}) bis 5 (\textit{stimme voll zu}) bewertet:

\begin{enumerate}
  \item Ich denke, dass ich diese App regelmäßig nutzen würde.
  \item Ich fand die App unnötig komplex.
  \item Ich fand die App einfach zu bedienen.
  \item Ich denke, ich würde die Unterstützung einer technischen Person
        benötigen, um diese App zu benutzen.
  \item Ich fand, dass die verschiedenen Funktionen gut integriert waren.
  \item Ich fand, dass es zu viele Inkonsistenzen in der App gab.
  \item Ich könnte mir vorstellen, dass die meisten Menschen die App
        schnell erlernen würden.
  \item Ich fand die App sehr umständlich zu bedienen.
  \item Ich fühlte mich sicher im Umgang mit der App.
  \item Ich musste viele Dinge lernen, bevor ich die App benutzen konnte.
\end{enumerate}

\subsubsection*{Abschlussinterview (offen)}

\begin{enumerate}
  \item Welche Funktion hat Ihnen am meisten gefallen? Warum?
  \item Welche Funktion hat Sie am meisten frustriert oder verwirrt?
  \item Würden Sie diese App nutzen, um Wien zu erkunden?
  \item Welche Funktion vermissen Sie?
  \item Haben Sie eine ähnliche App bereits verwendet?
        Was hat diese besser oder schlechter gemacht?
\end{enumerate}

\subsection{Durchführung}

Der Usability-Test wurde mit insgesamt \textbf{8 Testpersonen}
durchgeführt, je 2 Testpersonen pro Teammitglied.
Keine Testperson war Mitglied des Projektteams.
Jede Sitzung dauerte ca.\ 30–45 Minuten und folgte diesem Ablauf:

\begin{enumerate}
  \item Begrüßung und Einführung in die Think-Aloud-Methode (5~Min.)
  \item Durchführung der vier Aufgaben mit Post-Task-Fragen (20–30~Min.)
  \item Ausfüllen des SUS-Fragebogens (5~Min.)
  \item Abschlussinterview (5–10~Min.)
\end{enumerate}
